\documentclass[11pt,a4paper]{article}

\usepackage[utf8]{inputenc}
\usepackage[T1]{fontenc}
\usepackage[french]{babel}
\usepackage[a4paper,margin=2.2cm]{geometry}
\usepackage{amsmath,amssymb,amsthm,mathtools}
\usepackage{lmodern}
\usepackage{enumitem}
\usepackage{hyperref}
\usepackage{xcolor}
\usepackage{fancyhdr}
\usepackage{stmaryrd}
\usepackage{mathrsfs}
\usepackage{array}

\hypersetup{
    colorlinks=true,
    linkcolor=black,
    urlcolor=blue,
    citecolor=black
}

\setlength{\parindent}{0pt}
\setlength{\parskip}{0.6em}

\pagestyle{fancy}
\fancyhf{}
\fancyhead[L]{Algèbre linéaire --- Chapitre 12}
\fancyhead[R]{Espaces euclidiens}
\fancyfoot[C]{\thepage}

% Environnements
\newtheorem{definition}{Définition}[section]
\newtheorem{theoreme}[definition]{Théorème}
\newtheorem{proposition}[definition]{Proposition}
\newtheorem{corollaire}[definition]{Corollaire}
\newtheorem{lemme}[definition]{Lemme}
\theoremstyle{remark}
\newtheorem{remarque}[definition]{Remarque}
\newtheorem{exemple}[definition]{Exemple}
\newtheorem*{methode}{Méthode}
\newtheorem*{idee}{Idée}

% Commandes
\newcommand{\R}{\mathbb{R}}
\newcommand{\N}{\mathbb{N}}
\newcommand{\Z}{\mathbb{Z}}
\newcommand{\K}{\mathbb{K}}
\newcommand{\Vect}{\mathrm{Vect}}
\newcommand{\Ker}{\mathrm{Ker}}
\newcommand{\Img}{\mathrm{Im}}
\newcommand{\rg}{\mathrm{rg}}
\newcommand{\id}{\mathrm{Id}}
\newcommand{\tr}{\mathrm{tr}}
\newcommand{\Norm}[1]{\left\|#1\right\|}
\newcommand{\abs}[1]{\left|#1\right|}
\newcommand{\scal}[2]{\left\langle #1,#2\right\rangle}
\newcommand{\Orth}{\perp}
\newcommand{\Mnn}{M_n(\R)}
\newcommand{\On}{O_n(\R)}

\begin{document}

\begin{center}
    {\LARGE \textbf{Chapitre 12 --- Espaces euclidiens}}\\[0.4em]
    {\large Produit scalaire, orthogonalité et projections}
\end{center}

\hrule
\vspace{1em}

\tableofcontents

\newpage

\section{Pourquoi introduire les espaces euclidiens ?}

Jusqu’ici, on a étudié les espaces vectoriels sous un point de vue purement algébrique : combinaisons linéaires, bases, dimension, applications linéaires, réduction.

Mais dans beaucoup de situations, on veut aussi parler de :
\begin{itemize}
    \item longueur d’un vecteur ;
    \item angle entre deux vecteurs ;
    \item perpendicularité ;
    \item distance entre deux vecteurs ou entre un point et un sous-espace ;
    \item projections orthogonales.
\end{itemize}

Pour cela, la structure d’espace vectoriel ne suffit pas. Il faut introduire un nouvel outil : le \textbf{produit scalaire}.

\begin{idee}
Un espace euclidien est un espace vectoriel réel muni d’un produit scalaire, qui permet de faire de la géométrie dans un cadre algébrique.
\end{idee}

\section{Produit scalaire}

\subsection{Définition}

\begin{definition}
Soit \(E\) un espace vectoriel réel.  
Un \textbf{produit scalaire} sur \(E\) est une application
\[
\scal{\cdot}{\cdot}:E\times E\to \R
\]
vérifiant, pour tous \(x,y,z\in E\) et tout \(\lambda\in\R\) :
\begin{enumerate}[label=\arabic*)]
    \item \textbf{bilinéarité} :
    \[
    \scal{\lambda x+y}{z}=\lambda \scal{x}{z}+\scal{y}{z},
    \qquad
    \scal{x}{\lambda y+z}=\lambda \scal{x}{y}+\scal{x}{z} ;
    \]
    \item \textbf{symétrie} :
    \[
    \scal{x}{y}=\scal{y}{x} ;
    \]
    \item \textbf{positivité} :
    \[
    \scal{x}{x}\ge 0 ;
    \]
    \item \textbf{définie} :
    \[
    \scal{x}{x}=0 \Longleftrightarrow x=0.
    \]
\end{enumerate}
\end{definition}

\begin{definition}
Un \textbf{espace euclidien} est un espace vectoriel réel de dimension finie muni d’un produit scalaire.
\end{definition}

\begin{remarque}
On travaille ici sur \(\R\). Sur \(\C\), on parle d’espace hermitien, et le produit scalaire doit être sesquilinéaire.
\end{remarque}

\subsection{Exemple fondamental}

\begin{exemple}
Sur \(\R^n\), le produit scalaire canonique est défini par
\[
\scal{x}{y}=x_1y_1+\cdots+x_ny_n
\]
pour
\[
x=(x_1,\dots,x_n),\qquad y=(y_1,\dots,y_n).
\]
\end{exemple}

\begin{proof}
La bilinéarité et la symétrie sont immédiates. De plus,
\[
\scal{x}{x}=x_1^2+\cdots+x_n^2\ge 0,
\]
et
\[
\scal{x}{x}=0
\Longleftrightarrow
x_1=\cdots=x_n=0
\Longleftrightarrow
x=0.
\]
\end{proof}

\begin{exemple}
Sur \(\R_2[X]\), on peut définir
\[
\scal{P}{Q}=P(0)Q(0)+P(1)Q(1)+P(2)Q(2).
\]
C’est un produit scalaire.
\end{exemple}

\section{Norme et distance}

\subsection{Norme associée}

\begin{definition}
Dans un espace euclidien \((E,\scal{\cdot}{\cdot})\), on appelle \textbf{norme euclidienne} associée l’application
\[
\Norm{x}=\sqrt{\scal{x}{x}}.
\]
\end{definition}

\begin{remarque}
Le carré de la norme est
\[
\Norm{x}^2=\scal{x}{x}.
\]
\end{remarque}

\subsection{Distance}

\begin{definition}
La \textbf{distance} entre deux vecteurs \(x,y\in E\) est
\[
d(x,y)=\Norm{x-y}.
\]
\end{definition}

\begin{idee}
Dans un espace euclidien, on identifie souvent géométriquement les vecteurs à des points. La distance entre deux points est alors la norme de leur différence.
\end{idee}

\subsection{Premières propriétés}

\begin{proposition}
Pour tout \(x\in E\) et tout \(\lambda\in\R\),
\[
\Norm{x}\ge 0,
\qquad
\Norm{x}=0 \Longleftrightarrow x=0,
\qquad
\Norm{\lambda x}=\abs{\lambda}\Norm{x}.
\]
\end{proposition}

\begin{proof}
Les deux premières propriétés viennent directement de la définition du produit scalaire.

Pour la troisième :
\[
\Norm{\lambda x}^2=\scal{\lambda x}{\lambda x}=\lambda^2\scal{x}{x}=\lambda^2\Norm{x}^2.
\]
Comme les deux membres sont positifs,
\[
\Norm{\lambda x}=\abs{\lambda}\Norm{x}.
\]
\end{proof}

\section{Orthogonalité}

\subsection{Définition}

\begin{definition}
Deux vecteurs \(x,y\in E\) sont dits \textbf{orthogonaux} si
\[
\scal{x}{y}=0.
\]
On note alors
\[
x\Orth y.
\]
\end{definition}

\begin{exemple}
Dans \(\R^2\), les vecteurs \((1,0)\) et \((0,1)\) sont orthogonaux.
\end{exemple}

\begin{exemple}
Dans \(\R^3\), les vecteurs \((1,1,0)\) et \((1,-1,0)\) sont orthogonaux car
\[
1\cdot 1+1\cdot(-1)+0\cdot 0=0.
\]
\end{exemple}

\subsection{Orthogonal d’une partie}

\begin{definition}
Soit \(A\subset E\). On appelle \textbf{orthogonal} de \(A\), noté \(A^\Orth\), l’ensemble
\[
A^\Orth=\{x\in E\mid \forall a\in A,\ \scal{x}{a}=0\}.
\]
\end{definition}

\begin{proposition}
\(A^\Orth\) est un sous-espace vectoriel de \(E\).
\end{proposition}

\begin{proof}
Si \(x,y\in A^\Orth\) et \(\lambda,\mu\in\R\), alors pour tout \(a\in A\),
\[
\scal{\lambda x+\mu y}{a}
=
\lambda\scal{x}{a}+\mu\scal{y}{a}
=0.
\]
Donc \(\lambda x+\mu y\in A^\Orth\).
\end{proof}

\section{Inégalité de Cauchy-Schwarz}

\subsection{Énoncé}

\begin{theoreme}[Inégalité de Cauchy-Schwarz]
Pour tous \(x,y\in E\),
\[
\abs{\scal{x}{y}}\le \Norm{x}\Norm{y}.
\]
\end{theoreme}

\subsection{Démonstration}

\begin{proof}
Si \(y=0\), l’inégalité est évidente. Supposons \(y\neq 0\).

Pour tout \(t\in\R\), on a
\[
\Norm{x-ty}^2\ge 0.
\]
Développons :
\[
\Norm{x-ty}^2=\scal{x-ty}{x-ty}
=\Norm{x}^2-2t\scal{x}{y}+t^2\Norm{y}^2.
\]
C’est un polynôme du second degré en \(t\), toujours positif ou nul. Son discriminant est donc inférieur ou égal à zéro :
\[
\Delta=4\scal{x}{y}^2-4\Norm{x}^2\Norm{y}^2\le 0.
\]
Donc
\[
\scal{x}{y}^2\le \Norm{x}^2\Norm{y}^2,
\]
et en prenant la racine carrée :
\[
\abs{\scal{x}{y}}\le \Norm{x}\Norm{y}.
\]
\end{proof}

\subsection{Cas d’égalité}

\begin{proposition}
Il y a égalité dans Cauchy-Schwarz si et seulement si \(x\) et \(y\) sont colinéaires.
\end{proposition}

\begin{proof}
Dans la démonstration précédente, l’égalité a lieu si et seulement si le discriminant est nul, c’est-à-dire si et seulement si le polynôme
\[
\Norm{x-ty}^2
\]
admet une racine réelle \(t\), donc si et seulement si
\[
x-ty=0,
\]
c’est-à-dire si \(x\) et \(y\) sont colinéaires.
\end{proof}

\section{Inégalité triangulaire}

\begin{theoreme}
Pour tous \(x,y\in E\),
\[
\Norm{x+y}\le \Norm{x}+\Norm{y}.
\]
\end{theoreme}

\begin{proof}
On calcule :
\[
\Norm{x+y}^2=\Norm{x}^2+2\scal{x}{y}+\Norm{y}^2.
\]
Par Cauchy-Schwarz,
\[
2\scal{x}{y}\le 2\abs{\scal{x}{y}}\le 2\Norm{x}\Norm{y}.
\]
Donc
\[
\Norm{x+y}^2\le \Norm{x}^2+2\Norm{x}\Norm{y}+\Norm{y}^2=(\Norm{x}+\Norm{y})^2.
\]
Comme les deux membres sont positifs,
\[
\Norm{x+y}\le \Norm{x}+\Norm{y}.
\]
\end{proof}

\section{Identités remarquables}

\subsection{Polarisation}

\begin{proposition}[Identité de polarisation]
Pour tous \(x,y\in E\),
\[
\scal{x}{y}=\frac12\big(\Norm{x+y}^2-\Norm{x}^2-\Norm{y}^2\big).
\]
\end{proposition}

\begin{proof}
On développe :
\[
\Norm{x+y}^2=\Norm{x}^2+2\scal{x}{y}+\Norm{y}^2.
\]
On réarrange.
\end{proof}

\subsection{Parallélogramme}

\begin{proposition}
Pour tous \(x,y\in E\),
\[
\Norm{x+y}^2+\Norm{x-y}^2=2\Norm{x}^2+2\Norm{y}^2.
\]
\end{proposition}

\begin{proof}
On développe séparément :
\[
\Norm{x+y}^2=\Norm{x}^2+2\scal{x}{y}+\Norm{y}^2,
\]
\[
\Norm{x-y}^2=\Norm{x}^2-2\scal{x}{y}+\Norm{y}^2.
\]
En additionnant :
\[
\Norm{x+y}^2+\Norm{x-y}^2=2\Norm{x}^2+2\Norm{y}^2.
\]
\end{proof}

\begin{remarque}
C’est l’identité du parallélogramme.
\end{remarque}

\section{Familles orthogonales et orthonormales}

\subsection{Définitions}

\begin{definition}
Une famille \((e_1,\dots,e_p)\) est dite \textbf{orthogonale} si
\[
\forall i\neq j,\qquad \scal{e_i}{e_j}=0.
\]
Elle est dite \textbf{orthonormale} si, de plus,
\[
\Norm{e_i}=1
\qquad \text{pour tout } i.
\]
\end{definition}

\subsection{Liberté}

\begin{proposition}
Toute famille orthogonale de vecteurs non nuls est libre.
\end{proposition}

\begin{proof}
Supposons
\[
\lambda_1e_1+\cdots+\lambda_pe_p=0.
\]
En prenant le produit scalaire avec \(e_i\), on obtient
\[
\lambda_i\Norm{e_i}^2=0,
\]
car tous les autres termes sont nuls par orthogonalité. Comme \(\Norm{e_i}\neq 0\), on en déduit \(\lambda_i=0\). Ceci vaut pour tout \(i\).
\end{proof}

\begin{corollaire}
Dans un espace euclidien de dimension \(n\), une famille orthogonale de vecteurs non nuls contient au plus \(n\) vecteurs.
\end{corollaire}

\subsection{Coordonnées dans une base orthonormée}

\begin{theoreme}
Soit \((e_1,\dots,e_n)\) une base orthonormée de \(E\). Alors tout vecteur \(x\in E\) s’écrit
\[
x=\scal{x}{e_1}e_1+\cdots+\scal{x}{e_n}e_n.
\]
\end{theoreme}

\begin{proof}
Écrivons
\[
x=\lambda_1e_1+\cdots+\lambda_ne_n.
\]
En prenant le produit scalaire avec \(e_i\), on obtient
\[
\scal{x}{e_i}=\lambda_i,
\]
car la base est orthonormée.
\end{proof}

\begin{corollaire}[Parseval]
Dans une base orthonormée \((e_1,\dots,e_n)\),
\[
\Norm{x}^2=\scal{x}{e_1}^2+\cdots+\scal{x}{e_n}^2.
\]
\end{corollaire}

\begin{proof}
On remplace \(x\) par son développement dans la base orthonormée et on utilise l’orthonormalité.
\end{proof}

\section{Bases orthonormées}

\subsection{Existence}

\begin{theoreme}
Tout espace euclidien admet une base orthonormée.
\end{theoreme}

\begin{proof}
On part d’une base quelconque et on applique le procédé de Gram-Schmidt, puis on normalise les vecteurs obtenus.
\end{proof}

\subsection{Procédé de Gram-Schmidt}

\begin{theoreme}[Gram-Schmidt]
Soit \((u_1,\dots,u_n)\) une base d’un espace euclidien \(E\). On définit :
\[
v_1=u_1,
\]
puis, pour \(k\ge 2\),
\[
v_k=u_k-\sum_{j=1}^{k-1}\frac{\scal{u_k}{v_j}}{\Norm{v_j}^2}v_j.
\]
Alors \((v_1,\dots,v_n)\) est une famille orthogonale de vecteurs non nuls telle que
\[
\Vect(v_1,\dots,v_k)=\Vect(u_1,\dots,u_k)
\qquad \text{pour tout } k.
\]
\end{theoreme}

\begin{proof}
On raisonne par récurrence. Par construction, \(v_k\) est obtenu en retranchant à \(u_k\) sa projection sur l’espace engendré par \(v_1,\dots,v_{k-1}\). On vérifie que
\[
\scal{v_k}{v_i}=0
\quad \text{pour tout } i<k,
\]
puis que \(v_k\neq 0\) car sinon \(u_k\) appartiendrait à \(\Vect(u_1,\dots,u_{k-1})\), ce qui contredirait le fait que \((u_1,\dots,u_n)\) est une base.
\end{proof}

\begin{corollaire}
En posant
\[
e_i=\frac{v_i}{\Norm{v_i}},
\]
on obtient une base orthonormée \((e_1,\dots,e_n)\).
\end{corollaire}

\section{Sous-espaces orthogonaux}

\subsection{Orthogonal d’un sous-espace}

\begin{proposition}
Si \(F\) est un sous-espace vectoriel de \(E\), alors
\[
F^\Orth=\{x\in E\mid \forall y\in F,\ \scal{x}{y}=0\}
\]
est un sous-espace vectoriel de \(E\).
\end{proposition}

\begin{proof}
Déjà vu pour l’orthogonal d’une partie ; c’est en particulier vrai pour un sous-espace.
\end{proof}

\subsection{Somme directe orthogonale}

\begin{theoreme}
Si \(F\) est un sous-espace d’un espace euclidien \(E\), alors
\[
E=F\oplus F^\Orth.
\]
\end{theoreme}

\begin{proof}
Soit \((e_1,\dots,e_p)\) une base orthonormée de \(F\), que l’on complète en une base orthonormée
\[
(e_1,\dots,e_p,e_{p+1},\dots,e_n)
\]
de \(E\).

Alors
\[
F=\Vect(e_1,\dots,e_p),
\qquad
F^\Orth=\Vect(e_{p+1},\dots,e_n).
\]
La somme est donc directe et vaut \(E\).
\end{proof}

\begin{corollaire}
On a
\[
\dim(F)+\dim(F^\Orth)=\dim(E).
\]
\end{corollaire}

\begin{corollaire}
\[
(F^\Orth)^\Orth=F.
\]
\end{corollaire}

\begin{proof}
On a toujours
\[
F\subset (F^\Orth)^\Orth.
\]
Or
\[
\dim((F^\Orth)^\Orth)=\dim(E)-\dim(F^\Orth)=\dim(F).
\]
Donc l’inclusion est une égalité.
\end{proof}

\section{Projection orthogonale}

\subsection{Définition}

\begin{definition}
Soit \(F\) un sous-espace de \(E\). Pour tout \(x\in E\), comme
\[
E=F\oplus F^\Orth,
\]
il existe un unique couple \((p_F(x),q_F(x))\in F\times F^\Orth\) tel que
\[
x=p_F(x)+q_F(x).
\]
Le vecteur \(p_F(x)\) est appelé la \textbf{projection orthogonale de \(x\) sur \(F\)}.
\end{definition}

\subsection{Caractérisation}

\begin{proposition}
\(p_F(x)\) est l’unique vecteur de \(F\) tel que
\[
x-p_F(x)\in F^\Orth.
\]
\end{proposition}

\begin{proof}
C’est exactement la définition de la décomposition de \(x\) selon
\[
E=F\oplus F^\Orth.
\]
\end{proof}

\subsection{Expression dans une base orthonormée}

\begin{proposition}
Si \((e_1,\dots,e_p)\) est une base orthonormée de \(F\), alors
\[
p_F(x)=\sum_{i=1}^p \scal{x}{e_i}e_i.
\]
\end{proposition}

\begin{proof}
Le vecteur
\[
y=\sum_{i=1}^p \scal{x}{e_i}e_i
\]
appartient à \(F\). Pour tout \(j\),
\[
\scal{x-y}{e_j}=\scal{x}{e_j}-\sum_{i=1}^p \scal{x}{e_i}\scal{e_i}{e_j}=0.
\]
Donc \(x-y\in F^\Orth\), d’où \(y=p_F(x)\).
\end{proof}

\section{Distance à un sous-espace}

\subsection{Définition}

\begin{definition}
Soit \(F\) un sous-espace de \(E\), et \(x\in E\).  
La distance de \(x\) à \(F\) est
\[
d(x,F)=\inf_{y\in F}\Norm{x-y}.
\]
\end{definition}

\subsection{Théorème du projeté orthogonal}

\begin{theoreme}
Pour tout \(x\in E\),
\[
d(x,F)=\Norm{x-p_F(x)}.
\]
De plus, \(p_F(x)\) est l’unique vecteur de \(F\) qui réalise cette distance.
\end{theoreme}

\begin{proof}
Soit \(y\in F\). On écrit
\[
x-y=(x-p_F(x))+(p_F(x)-y).
\]
Or
\[
x-p_F(x)\in F^\Orth
\quad \text{et} \quad
p_F(x)-y\in F,
\]
donc ces deux vecteurs sont orthogonaux. Par Pythagore,
\[
\Norm{x-y}^2=\Norm{x-p_F(x)}^2+\Norm{p_F(x)-y}^2\ge \Norm{x-p_F(x)}^2.
\]
L’égalité a lieu si et seulement si
\[
p_F(x)-y=0,
\]
c’est-à-dire \(y=p_F(x)\).
\end{proof}

\begin{remarque}
La projection orthogonale est donc le meilleur approximant de \(x\) dans le sous-espace \(F\).
\end{remarque}

\section{Pythagore}

\begin{proposition}[Théorème de Pythagore]
Si \(x\Orth y\), alors
\[
\Norm{x+y}^2=\Norm{x}^2+\Norm{y}^2.
\]
\end{proposition}

\begin{proof}
On développe :
\[
\Norm{x+y}^2=\scal{x+y}{x+y}
=\Norm{x}^2+2\scal{x}{y}+\Norm{y}^2.
\]
Comme \(\scal{x}{y}=0\), on conclut.
\end{proof}

\section{Matrices du produit scalaire}

\subsection{Matrice dans une base}

\begin{definition}
Soit \(\mathcal B=(e_1,\dots,e_n)\) une base de \(E\). La matrice
\[
G=(\scal{e_i}{e_j})_{1\le i,j\le n}
\]
est appelée la \textbf{matrice de Gram} de la base \(\mathcal B\).
\end{definition}

\begin{proposition}
La matrice de Gram est symétrique.
\end{proposition}

\begin{proof}
Comme
\[
\scal{e_i}{e_j}=\scal{e_j}{e_i},
\]
on a
\[
G^T=G.
\]
\end{proof}

\begin{proposition}
Pour tout vecteur \(x=\sum x_ie_i\), on a
\[
\Norm{x}^2=X^T G X,
\]
où \(X\) est la colonne des coordonnées de \(x\) dans la base \(\mathcal B\).
\end{proposition}

\begin{proof}
On écrit
\[
x=\sum_i x_ie_i,
\]
puis
\[
\scal{x}{x}=\sum_{i,j}x_ix_j\scal{e_i}{e_j}.
\]
C’est exactement le produit matriciel \(X^TGX\).
\end{proof}

\subsection{Cas d’une base orthonormée}

\begin{proposition}
Une base \(\mathcal B\) est orthonormée si et seulement si sa matrice de Gram est \(I_n\).
\end{proposition}

\begin{proof}
Dans une base orthonormée,
\[
\scal{e_i}{e_j}=
\begin{cases}
1 & \text{si } i=j,\\
0 & \text{sinon}.
\end{cases}
\]
Donc \(G=I_n\). La réciproque est immédiate.
\end{proof}

\section{Matrices orthogonales}

\subsection{Définition}

\begin{definition}
Une matrice \(A\in M_n(\R)\) est dite \textbf{orthogonale} si
\[
A^TA=I_n.
\]
On note \(O_n(\R)\) l’ensemble des matrices orthogonales d’ordre \(n\).
\end{definition}

\begin{proposition}
Une matrice est orthogonale si et seulement si ses colonnes forment une base orthonormée de \(\R^n\).
\end{proposition}

\begin{proof}
Si \(C_1,\dots,C_n\) sont les colonnes de \(A\), alors le coefficient \((i,j)\) de \(A^TA\) est
\[
\scal{C_i}{C_j}.
\]
Ainsi
\[
A^TA=I_n
\]
équivaut à
\[
\scal{C_i}{C_j}=\delta_{ij},
\]
c’est-à-dire que les colonnes forment une base orthonormée.
\end{proof}

\begin{corollaire}
Si \(A\) est orthogonale, alors
\[
A^{-1}=A^T.
\]
\end{corollaire}

\begin{proof}
C’est précisément la relation
\[
A^TA=I_n.
\]
\end{proof}

\begin{corollaire}
Si \(A\) est orthogonale, alors
\[
\det(A)=\pm 1.
\]
\end{corollaire}

\begin{proof}
On prend le déterminant dans
\[
A^TA=I_n.
\]
On obtient
\[
\det(A^T)\det(A)=1.
\]
Comme \(\det(A^T)=\det(A)\),
\[
\det(A)^2=1.
\]
\end{proof}

\section{Isométries euclidiennes linéaires}

\subsection{Définition}

\begin{definition}
Un endomorphisme \(u:E\to E\) est une \textbf{isométrie euclidienne linéaire} si
\[
\forall x\in E,\qquad \Norm{u(x)}=\Norm{x}.
\]
\end{definition}

\subsection{Caractérisation par le produit scalaire}

\begin{proposition}
Un endomorphisme \(u\) est une isométrie euclidienne linéaire si et seulement si
\[
\forall x,y\in E,\qquad \scal{u(x)}{u(y)}=\scal{x}{y}.
\]
\end{proposition}

\begin{proof}
Si \(u\) conserve le produit scalaire, alors
\[
\Norm{u(x)}^2=\scal{u(x)}{u(x)}=\scal{x}{x}=\Norm{x}^2,
\]
donc \(u\) conserve la norme.

Réciproquement, si \(u\) conserve la norme, alors par l’identité de polarisation,
\[
\scal{u(x)}{u(y)}
=
\frac12\big(\Norm{u(x)+u(y)}^2-\Norm{u(x)}^2-\Norm{u(y)}^2\big).
\]
Comme \(u\) est linéaire,
\[
u(x)+u(y)=u(x+y),
\]
et comme elle conserve les normes,
\[
\scal{u(x)}{u(y)}=\frac12\big(\Norm{x+y}^2-\Norm{x}^2-\Norm{y}^2\big)=\scal{x}{y}.
\]
\end{proof}

\subsection{Version matricielle}

\begin{proposition}
Dans une base orthonormée, un endomorphisme est une isométrie euclidienne linéaire si et seulement si sa matrice est orthogonale.
\end{proposition}

\begin{proof}
Dans une base orthonormée, conserver le produit scalaire revient à conserver le produit scalaire canonique sur \(\R^n\), ce qui équivaut à
\[
A^TA=I_n.
\]
\end{proof}

\section{Applications classiques}

\subsection{Réflexion et projection}

\begin{proposition}
La projection orthogonale sur un sous-espace \(F\) est un projecteur :
\[
p_F^2=p_F.
\]
\end{proposition}

\begin{proof}
Si \(x=y+z\) avec \(y\in F\), \(z\in F^\Orth\), alors
\[
p_F(x)=y.
\]
Comme \(y\in F\),
\[
p_F(y)=y.
\]
Donc
\[
p_F(p_F(x))=p_F(y)=y=p_F(x).
\]
\end{proof}

\begin{proposition}
L’application
\[
s_F=2p_F-\id_E
\]
est la symétrie orthogonale par rapport à \(F\).
\end{proposition}

\begin{proof}
Si \(x=y+z\) avec \(y\in F\), \(z\in F^\Orth\), alors
\[
s_F(x)=2y-(y+z)=y-z.
\]
Cette transformation laisse \(F\) fixe et change le signe sur \(F^\Orth\).
\end{proof}

\section{Méthodes pratiques}

\begin{methode}
Pour montrer qu’une famille est orthogonale, on calcule tous les produits scalaires croisés.
\end{methode}

\begin{methode}
Pour montrer qu’une famille est orthonormale, on vérifie en plus que chaque vecteur est de norme \(1\).
\end{methode}

\begin{methode}
Pour calculer une projection orthogonale sur un sous-espace \(F\) :
\begin{enumerate}
    \item trouver une base orthonormée \((e_1,\dots,e_p)\) de \(F\) ;
    \item utiliser
    \[
    p_F(x)=\sum_{i=1}^p \scal{x}{e_i}e_i.
    \]
\end{enumerate}
\end{methode}

\begin{methode}
Pour obtenir une base orthonormée à partir d’une base quelconque, utiliser Gram-Schmidt puis normaliser.
\end{methode}

\section{Erreurs classiques}

\begin{itemize}
    \item Confondre orthogonalité et colinéarité.
    \item Oublier que le produit scalaire est défini positif.
    \item Croire qu’une famille orthogonale peut contenir le vecteur nul et rester une bonne base.
    \item Oublier de normaliser à la fin de Gram-Schmidt.
    \item Confondre projection orthogonale et simple décomposition dans une base.
    \item Se tromper entre \(F^\Orth\) et \((F^\Orth)^\Orth\).
    \item Oublier qu’une matrice orthogonale vérifie
    \[
    A^{-1}=A^T.
    \]
\end{itemize}

\section{Résumé du chapitre}

\begin{itemize}
    \item Un espace euclidien est un espace vectoriel réel de dimension finie muni d’un produit scalaire.
    \item La norme associée est
    \[
    \Norm{x}=\sqrt{\scal{x}{x}}.
    \]
    \item Deux vecteurs sont orthogonaux si leur produit scalaire est nul.
    \item L’inégalité de Cauchy-Schwarz donne
    \[
    \abs{\scal{x}{y}}\le \Norm{x}\Norm{y}.
    \]
    \item Toute famille orthogonale de vecteurs non nuls est libre.
    \item Toute base peut être orthonormalisée par le procédé de Gram-Schmidt.
    \item Pour tout sous-espace \(F\),
    \[
    E=F\oplus F^\Orth.
    \]
    \item La projection orthogonale sur \(F\) est le meilleur approximant dans \(F\).
    \item Dans une base orthonormée, les isométries euclidiennes linéaires sont représentées par les matrices orthogonales.
\end{itemize}

\section*{Exercices d’application immédiate}

\begin{enumerate}[label=\textbf{Exercice \arabic*.}]
    \item Vérifier que le produit canonique sur \(\R^n\) est un produit scalaire.

    \item Montrer que
    \[
    \scal{P}{Q}=P(0)Q(0)+P(1)Q(1)+P(2)Q(2)
    \]
    définit un produit scalaire sur \(\R_2[X]\).

    \item Calculer la norme des vecteurs
    \[
    (1,2),\qquad (1,-1,2),\qquad (3,0,4).
    \]

    \item Montrer que les vecteurs
    \[
    (1,1,0)\quad \text{et} \quad (1,-1,0)
    \]
    sont orthogonaux.

    \item Démontrer l’inégalité de Cauchy-Schwarz.

    \item Orthonormaliser la famille
    \[
    ((1,0),(1,1))
    \]
    dans \(\R^2\).

    \item Déterminer la projection orthogonale de \((1,2,3)\) sur la droite
    \[
    \Vect((1,1,1)).
    \]

    \item Déterminer la distance du vecteur \((1,2,3)\) au sous-espace
    \[
    \Vect((1,1,1)).
    \]

    \item Montrer qu’une matrice orthogonale vérifie
    \[
    A^{-1}=A^T.
    \]

    \item Soit \(F=\Vect((1,0,0),(0,1,0))\subset \R^3\). Déterminer \(F^\Orth\), puis la projection orthogonale sur \(F\).
\end{enumerate}

\end{document}